% part: intuitionistic-logic
% chapter: introduction
% section: constructive-reasoning

\documentclass[../../../include/open-logic-chapter]{subfiles}

\begin{document}

\olfileid[pt]{int}{int}{cr}

\section{Raciocínio construtivo}

Ao contrário das extensões da lógica clássica por operadores modais ou
quantificadores de segunda ordem, a lógica intuicionista é ``não
clássica'' porque restringe a lógica clássica. A lógica clássica é
\emph{não construtiva} de várias maneiras. A lógica intuicionista
pretende captar uma forma mais ``construtiva'' de raciocínio,
característica de certo tipo de matemática construtiva. Os exemplos a
seguir podem ilustrar algumas das motivações subjacentes.

Suponha que alguém afirmasse ter determinado um número natural~$n$ com
a propriedade de que, se $n$ é par, a hipótese de Riemann é verdadeira
e, se $n$ é ímpar, a hipótese de Riemann é falsa. Ótima notícia!{} Saber
se a hipótese de Riemann é verdadeira ou não é um dos grandes problemas
em aberto da matemática, e essa pessoa parece ter reduzido a questão a
um cálculo: determinar se certo número é par ou não.

Qual é o valor mágico de~$n$? A pessoa o descreve assim: $n$ é o número
natural igual a $2$ se a hipótese de Riemann for verdadeira, e igual a
$3$ caso contrário.

Com indignação, você exige seu dinheiro de volta. Do ponto de vista
clássico, a descrição acima de fato determina um único valor de $n$;
mas o que você realmente deseja é um valor de $n$ fornecido
\emph{explicitamente}.

Como outro exemplo, talvez menos artificial, considere a seguinte
questão. Sabemos que é possível elevar um número irracional a uma
potência racional e obter um resultado racional. Por exemplo,
$\sqrt{2}^2 = 2$. O que é menos claro é se é possível elevar um número
irracional a uma potência \emph{irracional} e obter um resultado
racional. O teorema a seguir responde afirmativamente:

\begin{thm}
Existem números irracionais $a$ e $b$ tais que $a^b$ é racional.
\end{thm}

\begin{proof}
Considere $\sqrt{2}^{\sqrt{2}}$. Se esse número é racional, terminamos:
podemos tomar $a = b = \sqrt{2}$. Caso contrário, ele é irracional.
Então temos
\[
(\sqrt{2}^{\sqrt{2}})^{\sqrt{2}} = \sqrt{2}^{\sqrt{2} \cdot
  \sqrt{2}} = \sqrt{2}^2 = 2,
\]
que é racional. Portanto, nesse caso, tome $a$ como
$\sqrt{2}^{\sqrt{2}}$ e $b$ como~$\sqrt 2$.
\end{proof}

Isso constitui uma prova válida? A maioria dos matemáticos considera
que sim. Mais uma vez, porém, há algo um pouco insatisfatório: provamos
a existência de um par de números reais com certa propriedade sem
poder dizer \emph{qual} é esse par. É possível provar o mesmo resultado
de modo que o par $a$, $b$ \emph{seja} fornecido na prova: tome
$a = \sqrt{3}$ e $b = \log_3 4$. Então
\[
a^b = \sqrt{3}^{\log_3 4} = 3^{1/2 \cdot \log_3 4} = (3^{\log_3
  4})^{1/2} = 4^{1/2}= 2,
\]
pois $3^{\log_3 x} = x$.

A lógica intuicionista é concebida para captar uma forma de raciocínio
na qual passos como o da primeira prova não são permitidos. Provar a
existência de um $x$ que satisfaça~$!A(x)$ significa ter de fornecer um
$x$ específico e uma prova de que ele satisfaz $!A$, como na segunda
prova. Provar que $!A$ ou $!B$ vale exige poder provar uma delas.

Formalmente, a lógica intuicionista é o que se obtém ao restringir de
certo modo um sistema de !!a{derivation} da lógica clássica. Do ponto
de vista matemático, trata-se apenas de sistemas dedutivos formais;
contudo, como já observamos, eles pretendem captar uma forma de
raciocínio matemático. Pode-se entendê-la como a forma de raciocínio
justificada por certa visão filosófica da matemática (como o
intuicionismo de Brouwer); pode-se entendê-la como uma forma mais
``concreta'' e satisfatória de raciocínio matemático (na linha do
construtivismo de Bishop); e pode-se discutir se a descrição formal
capta ou não a motivação informal. Quaisquer que sejam, porém, nossas
posições filosóficas, podemos estudar a lógica intuicionista como uma
lógica apresentada formalmente; e, por motivos diversos, muitos
lógicos matemáticos consideram interessante fazê-lo.

\end{document}
